%!TEX root = ../mainufficiale.tex
\chapter{Esperimenti e risultati}
\label{cap:confronto}

Questo capitolo descrive il protocollo sperimentale di questa tesi e ne
riporta i risultati. Definisce anzitutto le metriche di valutazione
(§\ref{sec:metriche}) e il protocollo sperimentale
(§\ref{sec:protocollo}); presenta poi il confronto di base tra le
baseline classiche e MetaSTGAT Pro insieme ai tre studi sperimentali di
questa tesi, sul peso $\alpha$ della componente di equità, sull'ablation
della procedura di addestramento e sui meccanismi di comunicazione
spaziale alternativi (§\ref{sec:confronti}), e ne riporta i risultati
(§\ref{sec:risultati-baseline}).

\section{Metriche di valutazione}
\label{sec:metriche}

Nella letteratura sul Traffic Signal Control le metriche di valutazione
proposte sono numerose.
Questo capitolo introduce le 4 metriche principali utilizzate, coerentemente con il duplice obiettivo di
questo lavoro, ovvero efficienza ed equità (come motivato nel Capitolo~\ref{cap:introduzione}):
il \emph{travel time} (tempo di viaggio, medio e massimo) e il \emph{trip completion} (veicoli che
completano il percorso) per l'efficienza, il \emph{max wait} (l'attesa
massima registrata da un veicolo ad un incrocio) per l'equità.
C'è da sottolineare che né \emph{travel time} né
\emph{trip completion} sono metriche che appaiono esplicitamente nella
ricompensa dei modelli. Pertanto qualunque risultato su di esse sarà
l'indiretta conseguenza di come la reward ha indirizzato l'apprendimento
della policy.

\paragraph{Travel time.} Per ogni veicolo $v$, il tempo di viaggio è
\begin{equation}
  tt_v =
  \begin{cases}
    t_{\text{arrivo}}(v) - t_{\text{spawn}}(v)     & \text{se } v \text{ è arrivato entro la fine dell'episodio} \\
    t_{\text{fine episodio}} - t_{\text{spawn}}(v) & \text{altrimenti (un lower bound del vero valore)}
  \end{cases}
\end{equation}
Una delle due metriche cardine in questo studio è il travel time medio $\overline{tt}$,
ovvero la media di $ tt_v $ calcolata sull'intera popolazione di veicoli generati nell'episodio.
Tale metrica non considera soltanto i veicoli arrivati, ma include anche i
veicoli che non hanno completato il percorso entro la fine della simulazione.
La ragione è evitare un bias di sopravvivenza: un
modello che congestiona la rete e lascia arrivare solo una minoranza di
veicoli su percorsi ottimizzati otterrebbe, calcolando la media solo sui
veicoli arrivati, un risultato positivo ma non rappresentativo del comportamento reale della rete.

La seconda vera metrica cardine in questo lavoro di tesi è il travel time massimo $\overline{tt}_{\max}$,
ovvero il massimo tra tutti i travel time dei veicoli generati, anche in questo caso indipendentemente
dal completamento del percorso, per motivi analoghi a quelli descritti per il travel time medio.

\paragraph{Trip completion.} Questa metrica misura la percentuale di veicoli che hanno completato il percorso entro la
fine dell'episodio, escludendo dal totale i veicoli strutturalmente impossibilitati a farlo: quelli il cui tragitto, percorso a velocità massima costante e senza alcun arresto ai semafori, richiederebbe già da solo più dei secondi restanti alla fine dell'episodio a partire dal proprio istante di generazione.
Questo sottoinsieme dipende solo dallo scenario di traffico (percorso e istante di generazione di ciascun veicolo), non dal controllore, e viene quindi calcolato una sola volta per configurazione e riusato per ogni modello.
In questo studio questa metrica viene indicata con $tc$. Il dettaglio di
$tc$ per ogni singolo modello e ogni singola configurazione è in
Tabella~\ref{tab:tc-adjusted}, a fine capitolo.

\paragraph{Max wait.} Questa metrica riguarda il tempo più lungo che un
veicolo trascorre su una singola corsia del proprio percorso. Sia $w_v$
questo tempo per il veicolo $v$: l'intervallo tra l'istante in cui $v$
entra in una corsia del suo tragitto e quello in cui supera l'incrocio a
cui essa conduce, calcolato come il massimo su tutte le corsie
attraversate lungo il percorso. Il conteggio su una corsia non si azzera
se il veicolo avanza al suo interno, e prosegue anche se il veicolo si è
fermato e ripartito una o più volte prima di superare l'incrocio; si
interrompe soltanto quando la vettura lo supera. Il max wait
dell'episodio è
\begin{equation}
  \text{wait}_{\max} = \max_{v}\ w_v ,
\end{equation}
il massimo su tutti i veicoli generati nell'episodio.

Come per il travel time, include anche gli stalli ancora in corso al
termine dell'episodio. Sulle configurazioni con spaziatura di 200 metri
(§\ref{subsec:dataset}), il conteggio è esteso all'intera corsia invece di
limitarsi al campo visivo $d_{\text{vis}}$ dei modelli di reinforcement
learning (§\ref{sec:formalizzazione}): un veicolo fermo oltre $d_{\text{vis}}$
dal semaforo sarebbe altrimenti escluso dal computo, sottostimando l'attesa
peggiore effettivamente osservata sulla corsia.
In questo studio max wait viene indicato con $\text{wait}_{\max}$.

La presente metrica non va confusa con $\mathrm{MaxRedWait}_i^t$, il termine
usato nella ricompensa (Capitolo~\ref{cap:tsc-modelli}) per
penalizzare l'attesa durante il training. $\mathrm{MaxRedWait}_i^t$ è
calcolato a ogni singolo passo di decisione, su una sola intersezione alla
volta, ma soprattutto è limitata al campo visivo dell'incrocio.
$\text{wait}_{\max}$ è invece calcolato una sola volta a fine
episodio, come massimo su tutta la rete, e fa riferimento all'intera corsia.

\section{Protocollo sperimentale}
\label{sec:protocollo}

Ogni modello RL è addestrato sulle tre configurazioni di training presentate nella Tabella~\ref{tab:split}.

Seguendo un ciclo fisso, le tre configurazioni si alternano episodio per
episodio. A fine training, l'ultimo checkpoint e quello che ha ottenuto il
valore di loss più basso durante l'intero addestramento sono confrontati
sulla configurazione di validazione \texttt{config\_4x4\_100m\_6k\_flat},
mediando $\overline{tt}$ su 5 seed diversi.
Questa fase serve a evitare due problemi che si concatenano tra loro.
Il primo è che in un ciclo di
apprendimento online la politica non converge necessariamente in modo
monotono, dunque gli ultimi episodi addestrati potrebbero allontanare la policy da un
punto di minimo quantomeno locale.
Il secondo è che la scelta del modello che ha registrato la loss minima non garantisce
le capacità di generalizzazione del modello, in quanto una loss bassa potrebbe
anche essere sintomo di una policy che ha soltanto memorizzato delle condizioni particolari delle
configurazioni di training, senza però aver realmente appreso una politica in grado di generalizzare.

La selezione si basa sul solo $\overline{tt}$ e non sulle metriche di caso
peggiore che questo lavoro tiene in primo piano altrove: mediare il
travel time medio sui 5 seed di validazione dà già una stima
ragionevolmente stabile del checkpoint migliore, mentre fare lo stesso con
un massimo come $tt_{\max}$ o $\text{wait}_{\max}$ resterebbe una stima
molto più rumorosa, perché un massimo risente della coda della
distribuzione assai più di una media. È un limite pratico del protocollo,
non una scelta di principio.

Il modello che risulta essere migliore in questa fase di selezione, viene poi utilizzato nella fase di test.
Ogni valutazione, sia in validation sia nella batteria di test che descriveremo ora, avviene con $\varepsilon=0$:
nessuna esplorazione residua, solo la politica appresa allo stato puro. Inoltre i risultati ottenuti in ciascuna
configurazione testata, sono mediati su 5 seed diversi, analogamente a quanto fatto per $\overline{tt}$ nella fase di validazione.

Misurare la capacità di generalizzazione dei modelli è uno degli obiettivi
dichiarati di questo lavoro (Capitolo~\ref{cap:introduzione}), un asse che
gran parte della letteratura sul controllo semaforico appreso non copre
sistematicamente, in quanto la prassi più diffusa, seguita anche dall'articolo
originale di MetaSTGAT \citep{wang2022metastgat}, è addestrare e valutare un modello sulla stessa
rete stradale, senza mai trasferirlo a una topologia diversa.

Per questo motivo la batteria di test di questo lavoro include sia le 3
configurazioni di training sia 7 configurazioni aggiuntive mai viste
durante l'addestramento. Le prime permettono un confronto diretto con la
prassi di letteratura appena descritta, riportando un risultato ottenuto
nelle stesse condizioni con cui la maggior parte degli altri lavori
riporta il proprio; le seconde, assenti nella maggior parte degli studi
citati, misurano invece se le prestazioni osservate si confermano anche con il cambio
di topologia e di condizioni di traffico.

Ogni modello, RL o baseline, è dunque valutato sistematicamente sulle
stesse 10 configurazioni: le 3 varianti di training stesse (per misurare la
prestazione in-distribuzione) più 7 configurazioni mai viste in training
(4x4 a 100m con profilo \emph{peak}, 4x4 a 200m invece di 100m con densità sia
\emph{flat} sia \emph{peak}, $5\times5$ e $6\times6$ a densità sia
\emph{flat} sia \emph{peak}), usate per misurare la generalizzazione a
profili di traffico, spaziature e topologie mai incontrati in training.

Ogni modello di questo lavoro, RL o ablation, deriva da un singolo run di
addestramento. La batteria a più seed appena descritta ripete la sola
valutazione, non l'addestramento: le medie e le deviazioni standard
riportate in questo capitolo descrivono quindi la variabilità della
policy già addestrata su scenari di valutazione diversi, non la
variabilità tra run di addestramento indipendenti.

\section{Confronti proposti}
\label{sec:confronti}

Il confronto sperimentale di questo capitolo è organizzato in tre studi indipendenti,
ciascuno dei quali varia una sola dimensione (il peso $\alpha$, un gruppo di
modifiche alla procedura, il meccanismo spaziale) rispetto al
modello di riferimento ($\alpha=0.08$, MetaSTGAT a un layer) a cui ci riferiremo con il termine \emph{MetaSTGAT Pro},
lasciando invariato tutto il resto. Verranno analizzati nel seguito i tre studi proposti.
Ciascun confronto conterrà al suo interno, come riferimento, le due baseline classiche, Fixed-Time e Max-Pressure, e il modello \emph{MetaSTGAT Pro}.

\subsection{Studio del peso \texorpdfstring{$\alpha$}{alpha}}

Il primo studio (§\ref{sec:risultati-alpha}) pone a confronto tre varianti del
modello MetaSTGAT Pro che differiscono solo per il valore del peso $\alpha$ riguardante la penalità anti-starvation
nella ricompensa (Eq.~\ref{eq:reward-custom}).
In particolare vengono valutati i valori di $\alpha \in \{0.08, 0.04, 0.00\}$.
È il confronto più diretto per quantificare il compromesso tra efficienza ed
equità dichiarato come obiettivo di questo lavoro.
Variando solo questo iperparametro della ricompensa, senza toccare stato o architettura,
osserveremo quanto costa in termini di travel time medio ottenere una data riduzione del travel time massimo e del max wait.

\subsection{Studio di ablazione}

Il secondo confronto (§\ref{sec:ablation-risultati}) isola sistematicamente ogni blocco di differenze
tra l'implementazione di questo lavoro e la procedura
descritta nell'articolo originale di MetaSTGAT, tramite 6 preset che
vanno dalla replica fedele del paper al modello MetaSTGAT Pro.
È un confronto necessario perché né l'uno né l'altro estremo, da soli,
spiegano il divario di prestazioni osservato tra i due: isolare un
sottosistema alla volta permette di valutare il contributo di ciascuna modifica.

Ogni preset applica un gruppo coerente di scelte configurative,
isolando un sottosistema alla volta rispetto al modello MetaSTGAT Pro (§\ref{subsec:framing-pro}):

\begin{table}[htbp]
  \centering
  \small
  \begin{tabular}{@{}llp{8.5cm}@{}}
    \toprule
    Preset                   & Reward          & Cosa isola rispetto a Pro                                            \\
    \midrule
    \texttt{pro}             & \texttt{custom} & --                                                                   \\
    \texttt{temporal}        & \texttt{custom} & burn-in + BPTT sulla Meta-LSTM                                       \\
    \texttt{single\_dqn}     & \texttt{custom} & Double DQN                                                           \\
    \texttt{vanilla\_buffer} & \texttt{custom} & PER, Huber, grad clipping, warm-up, esplorazione ciclica             \\
    \texttt{environment}     & \texttt{paper}  & campo visivo, feature dello stato $w_i^t$, maschera fasi consecutive \\
    \texttt{paper}           & \texttt{paper}  & Replica fedele di Wang et al. (2022)                                 \\
    \bottomrule
  \end{tabular}
  \caption{Preset di ablazione e sottosistema isolato da ciascuno. La reward \texttt{custom} è definita da Eq. (\ref{eq:reward-custom}), mentre la reward \texttt{paper} è definita da Eq. (\ref{eq:reward-paper}).}
  \label{tab:ablation-preset}
\end{table}

A differenza degli altri quattro preset, che modificano un solo sottosistema,
\texttt{environment} cambia contemporaneamente osservazione (campo visivo,
$w_i^t$, maschera) e ricompensa (torna a quella di Wang et al., 2022): i
suoi risultati misurano quindi l'effetto congiunto delle due cose, non
isolano un singolo meccanismo.

\subsection{Studio sui meccanismi spaziali}

Il terzo studio (§\ref{sec:confronto-spaziale}) sostituisce il solo
meccanismo di comunicazione spaziale tra intersezioni di MetaSTGAT Pro.
Vengono testate le seguenti varianti già introdotte nel Capitolo \ref{sec:spatial-alt}:
\begin{itemize}
  \item MetaSTGAT (il riferimento, a un layer come nella formulazione standard e a due layer)
  \item MetaSTGCN (a uno e due layer)
  \item MetaSTSONAR (con $L=2$ e $L=4$ iterazioni)
\end{itemize}

È il confronto che verifica se il beneficio osservato dipenda
specificamente dal meccanismo di attenzione appreso, o se una qualunque
forma di condivisione spaziale dell'informazione tra nodi vicini produca
risultati equivalenti.

\section{Risultati}
\label{sec:risultati-baseline}

I tre studi del precedente capitolo (§\ref{sec:confronti}) sono presentati
secondo uno schema comune: una tabella mediata sulle 3 configurazioni di
training, una corrispondente sulle 7 di test, un grafico a bande sovrapposte che le
riassume visivamente, un plot della variazione percentuale delle metriche principali rispetto a
Max-Pressure, gli eventuali casi particolari osservabili solo scendendo
alla singola configurazione, chiusi da una discussione che include, dove il risultato non
era quello atteso, una possibile spiegazione.
In queste tabelle, dallo studio sul peso $\alpha$ in poi, sono evidenziati
in grassetto i due risultati migliori di ogni colonna; nelle due tabelle
del confronto di base (§\ref{sec:risultati-base}), che hanno solo tre
righe, resta invece in grassetto il solo valore migliore.

Fixed-Time, Max-Pressure e MetaSTGAT Pro ($\alpha=0.08$, il modello di
riferimento di questo lavoro) compaiono in ciascuno dei tre studi.
Dunque, per non ripetere la stessa discussione su questi tre riferimenti di base,
la sottosezione seguente li confronta da soli,
mentre le sottosezioni successive approfondiranno i rispettivi confronti
(peso $\alpha$, ablazione, meccanismi spaziali).

\subsection{Fixed-Time, Max-Pressure e Pro: il confronto di base}
\label{sec:risultati-base}

\begin{table}[htbp]
  \centering
  \begin{tabular}{@{}lrrrr@{}}
    \toprule
    Modello                       & $\overline{tt}$ (s) & $tt_{\max}$ (s) & wait$_{\max}$ (s) & $tc$ $(\%)$   \\
    \midrule
    Fixed-Time                    & 339.3               & 1030.0          & 935.0             & 86.2          \\
    Max-Pressure                  & \textbf{161.9}      & 1070.0          & 1040.0            & \textbf{97.9} \\
    MetaSTGAT Pro ($\alpha=0.08$) & 172.8               & \textbf{630.0}  & \textbf{425.0}    & 97.4          \\
    \bottomrule
  \end{tabular}
  \caption{Travel time medio, massimo, max wait e trip completion dei tre
    controllori presenti in ogni studio di questo capitolo, mediati sulle 3
    configurazioni di training.}
  \label{tab:base-train}
\end{table}

\begin{table}[htbp]
  \centering
  \begin{tabular}{@{}lrrrr@{}}
    \toprule
    Modello                       & $\overline{tt}$ (s) & $tt_{\max}$ (s) & wait$_{\max}$ (s) & $tc$ $(\%)$   \\
    \midrule
    Fixed-Time                    & 460.7               & 1496.1          & 1311.0            & 84.0          \\
    Max-Pressure                  & \textbf{238.3}      & 1647.4          & 1636.7            & 91.9          \\
    MetaSTGAT Pro ($\alpha=0.08$) & 241.0               & \textbf{1017.4} & \textbf{792.0}    & \textbf{93.8} \\
    \bottomrule
  \end{tabular}
  \caption{Come Tabella~\ref{tab:base-train}, mediata sulle 7 configurazioni
    di test.}
  \label{tab:base-test}
\end{table}

Fixed-Time è il controllore peggiore su travel time medio e su trip
completion, in entrambi i regimi. Il quadro si inverte però su
$tt_{\max}$ e $wait_{\max}$: Max-Pressure non batte Fixed-Time su nessuna
delle due metriche, né in training ($tt_{\max}$ 1070.0s contro 1030.0s,
$wait_{\max}$ 1040.0s contro 935.0s) né in test (1647.4s contro 1496.1s,
1636.7s contro 1311.0s). Una spiegazione plausibile è strutturale:
Max-Pressure sceglie a ogni istante la fase che libera più veicoli, senza
mai considerare da quanto tempo attende la direzione esclusa
(Capitolo~\ref{cap:introduzione}), e può quindi in linea di principio far
attendere a lungo una direzione a bassa pressione se le altre restano
sempre più cariche; la rotazione fissa di Fixed-Time, per quanto cieca al
traffico reale, garantisce invece per costruzione che ogni fase torni
periodicamente, ponendo un limite strutturale all'attesa massima che una
regola puramente reattiva non ha. MetaSTGAT Pro resta l'unico dei tre a
battere Fixed-Time su ogni metrica, in entrambi i regimi. Non osservando
mai lo stato della rete, Fixed-Time resta comunque il riferimento più
semplice rispetto a cui ogni strategia reattiva o appresa può essere
confrontata.

Tra Max-Pressure e MetaSTGAT Pro il confronto è più sfumato. In training e in test Max-Pressure
resta il modello con il travel time medio più basso in assoluto (rispettivamente 161.9s e 238.3s),
coerente con la sua ottimalità di portata a livello di rete;
MetaSTGAT Pro paga per questo un travel time leggermente peggiore sia in training (172.8s, +6.7\%) sia in test (241.0s, +1.1\%).
Un'eccezione a questo quadro emerge scendendo alle singole configurazioni
di test: sulle due griglie $6\times6$, le più estese testate e mai viste
in training, MetaSTGAT Pro registra un travel time medio inferiore a
quello di Max-Pressure (238.1s contro 244.7s in \emph{flat}, 275.1s contro
283.3s in \emph{peak}), l'unico caso tra le 10 configurazioni in cui questo
avviene. Lo stesso pattern, sulle stesse due configurazioni, si osserva
anche per $\alpha=0.04$ e per \texttt{vanilla\_buffer}
(§\ref{sec:ablation-risultati}), il che rende meno probabile si tratti di
rumore. Una spiegazione plausibile è che, sulla rete più estesa e mai
vista in training, il coordinamento appreso tra intersezioni inizi a
incidere anche sull'efficienza media, mentre Max-Pressure, privo di
qualunque nozione della rete oltre l'incrocio corrente, perda parte del
proprio vantaggio quando la rete cresce oltre la scala di training; questo
lavoro non arriva a verificare se il divario continui a crescere su reti
ancora più estese.
Per quanto riguarda il $wait_{\max}$ il rapporto si inverte nettamente:
MetaSTGAT Pro registra un valore il 59\% più basso di Max-Pressure in training e il 52\% più basso in test
(Tabelle~\ref{tab:base-train}-\ref{tab:base-test}).
Questi risultati introducono il compromesso al centro di questo lavoro, discusso nella sua forma
completa al variare di $\alpha$ nella prossima sottosezione.
Il vantaggio relativo di MetaSTGAT Pro sul $wait_{\max}$ di Max-Pressure è
più ampio in training (59\%) che in test (52\%); una lettura cauta è che
parte di questo vantaggio si eroda, come atteso, su condizioni mai
incontrate in addestramento, senza che questo permetta di per sé
conclusioni più forti sulle capacità di generalizzazione del modello.

Sul completamento ($tc$) la separazione più netta è quella tra Fixed-Time
e tutti gli altri modelli di questo capitolo: 86.2\%/84.0\%
(training/test) contro un intervallo compreso tra il 94.4\% e il 97.9\% in
training e tra l'89.0\% e il 95.8\% in test per ogni altra variante
(Tabella~\ref{tab:tc-adjusted}, a fine capitolo); l'unica
eccezione parziale è MetaSTGCN, il cui completamento più basso è discusso
a parte in §\ref{sec:confronto-spaziale}. Tra Max-Pressure e MetaSTGAT
Pro, quest'ultimo ottiene un $tc$ più alto su tutte e 7 le configurazioni
di test, un vantaggio sistematico non limitato alle due griglie
$6\times6$ già osservate sul travel time medio. Con questo unico scarto
rilevante già osservato qui, le sottosezioni seguenti non ripetono la
discussione di $tc$ configurazione per configurazione: si limitano a
segnalarne la tendenza generale, rimandando alla stessa tabella per il
dettaglio completo.

\subsection{Studio del peso \texorpdfstring{$\alpha$}{alpha}: risultati}
\label{sec:risultati-alpha}
\label{sec:equita-alpha}

Questa sottosezione confronta i due controllori di riferimento
(§\ref{sec:risultati-base}) con le tre varianti del modello Pro che
differiscono solo per il peso $\alpha$ della penalità anti-starvation,
con $\alpha \in \{0.08, 0.04, 0.00\}$.

\begin{table}[htbp]
  \centering
  \begin{tabular}{@{}lrrr@{}}
    \toprule
    Modello                       & $\overline{tt}$ (s) & $tt_{\max}$ (s) & wait$_{\max}$ (s) \\
    \midrule
    Fixed-Time                    & 339.3               & 1030.0          & 935.0             \\
    Max-Pressure                  & \textbf{161.9}      & 1070.0          & 1040.0            \\
    MetaSTGAT Pro ($\alpha=0.08$) & 172.8               & \textbf{630.0}  & \textbf{425.0}    \\
    MetaSTGAT Pro ($\alpha=0.04$) & \textbf{170.7}      & \textbf{685.0}  & 560.0             \\
    MetaSTGAT Pro ($\alpha=0.00$) & 173.5               & 740.0           & \textbf{555.0}    \\
    \bottomrule
  \end{tabular}
  \caption{Travel time medio, massimo e max wait, mediati
    sulle 3 configurazioni di training.}
  \label{tab:risultati-train}
\end{table}

\begin{table}[htbp]
  \centering
  \begin{tabular}{@{}lrrr@{}}
    \toprule
    Modello                       & $\overline{tt}$ (s) & $tt_{\max}$ (s) & wait$_{\max}$ (s) \\
    \midrule
    Fixed-Time                    & 460.7               & 1496.1          & 1311.0            \\
    Max-Pressure                  & \textbf{238.3}      & 1647.4          & 1636.7            \\
    MetaSTGAT Pro ($\alpha=0.08$) & \textbf{241.0}      & \textbf{1017.4} & \textbf{792.0}    \\
    MetaSTGAT Pro ($\alpha=0.04$) & 241.1               & \textbf{1054.7} & \textbf{930.9}    \\
    MetaSTGAT Pro ($\alpha=0.00$) & 247.5               & 1164.4          & 1012.3            \\
    \bottomrule
  \end{tabular}
  \caption{Come Tabella~\ref{tab:risultati-train}, mediata sulle 7
    configurazioni di test.}
  \label{tab:risultati-test}
\end{table}

Sul travel time medio ciò che verrebbe naturale attendersi è che
al crescere del valore di $\alpha$ la metrica in questione peggiori, con
le due configurazioni prive di penalità anti-starvation, Max-Pressure e $\alpha=0.00$, che risulterebbero essere le
migliori, in quanto nessuna delle due paga un compromesso con l'equità, quindi
nessuna dovrebbe peggiorare travel time medio in cambio di un vantaggio che non
persegue.
I dati non confermano questa attesa in modo netto: $\alpha=0.00$ non è la
variante Pro con il travel time medio più basso, né in training né in
test (173.5s in training contro 172.8s di $\alpha=0.08$ e 170.7s di
$\alpha=0.04$; 247.5s in test contro 241.0s e 241.1s), ma il distacco dalle
altre due varianti non è sistematico come atteso. Una
causa plausibile: con
$\alpha=0$ il termine di ricompensa legato all'attesa scompare, ma il
vettore $w_i^t$ resta nello stato osservato dalla rete
(§\ref{sec:formalizzazione}); la rete continua quindi a vedere un segnale
che la ricompensa non penalizza più, un disallineamento tra stato e
obiettivo che può disturbare l'addestramento.

Tra le tre varianti Pro, $\alpha=0.04$ resta quella con il travel time
medio più basso in training (170.7s), pur restando dietro a Max-Pressure
in entrambi i regimi.

\begin{figure}[htbp]
  \centering
  \begin{subfigure}[t]{0.48\textwidth}
    \centering
    \includegraphics[width=\linewidth]{pro_alpha_train.png}
    \caption{Configurazioni di training}
  \end{subfigure}
  \hfill
  \begin{subfigure}[t]{0.48\textwidth}
    \centering
    \includegraphics[width=\linewidth]{pro_alpha_test.png}
    \caption{Configurazioni di test}
  \end{subfigure}
  \caption{Travel time medio (fascia scura), max wait (fascia intermedia) e
    travel time massimo (fascia chiara) per lo studio sistematico su $\alpha$, stessi dati
    delle Tabelle~\ref{tab:risultati-train}-\ref{tab:risultati-test}.}
  \label{fig:overlay-pro-alpha}
\end{figure}

\begin{figure}[htbp]
  \centering
  \begin{subfigure}[t]{0.48\textwidth}
    \centering
    \includegraphics[width=\linewidth]{diff_pct_alpha_train.png}
    \caption{Configurazioni di training}
  \end{subfigure}
  \hfill
  \begin{subfigure}[t]{0.48\textwidth}
    \centering
    \includegraphics[width=\linewidth]{diff_pct_alpha_test.png}
    \caption{Configurazioni di test}
  \end{subfigure}
  \caption{Variazione percentuale di travel time medio e massimo rispetto a
    Max-Pressure (riferimento a 0\%) per lo studio su $\alpha$, calcolata per
    singola configurazione e poi mediata.}
  \label{fig:diffpct-alpha}
\end{figure}

Su $tt_{\max}$ il quadro è quello atteso: la metrica peggiora in modo
strettamente monotono al diminuire di $\alpha$, cioè al ridursi del peso
dato al caso di starvation più estremo, in entrambi i regimi. Su
$wait_{\max}$ l'andamento è quasi monotono: in test $\alpha=0.00$ (1012.3s)
resta peggiore di $\alpha=0.04$ (930.9s) come atteso, ma in training i due
valori sono sostanzialmente pari (555.0s contro 560.0s): le configurazioni
di training, a differenza di quelle di test, non sono mediate su più
seed, per cui una differenza di 5 secondi tra le due non è distinguibile
da una singola realizzazione senza una stima della variabilità a
supporto. Il
riferimento $\alpha=0.08$ ottiene comunque sia il $tt_{\max}$ sia il
$wait_{\max}$ più bassi in entrambi i regimi, il 59\% in meno
di Max-Pressure in training e il 52\% in meno in test. In test sia $\alpha=0.04$
sia $\alpha=0.00$ restano nettamente sotto Fixed-Time su entrambe le metriche
($wait_{\max}$ rispettivamente 930.9s e 1012.3s, $tt_{\max}$ 1054.7s e 1164.4s,
contro 1311.0s e 1496.1s di Fixed-Time): rimuovere del tutto la
penalità anti-starvation non avvicina le prestazioni di caso peggiore
a quelle di un controllo non adattivo.

Nonostante il lieve vantaggio di $\alpha=0.04$ sul travel time medio in
training, il resto di questo lavoro adotta $\alpha=0.08$ come riferimento
standard per gli studi di ablazione e dei meccanismi spaziali
(§\ref{sec:ablation-risultati}, §\ref{sec:confronto-spaziale}). La ragione
è la priorità dichiarata tra le metriche di questo lavoro:
$\alpha=0.04$ è sistematicamente meno incisivo di
$\alpha=0.08$ su $wait_{\max}$ e $tt_{\max}$, le due metriche di caso peggiore
che questa tesi tiene in primo piano, a fronte di un guadagno sul solo
travel time medio, la meno critica delle tre e comunque assente in test.
$\alpha=0.08$ resta quindi
la scelta più coerente con l'obiettivo del lavoro, anche al costo di
qualche secondo in più di travel time medio.

\subsection{Studio di ablazione: risultati}
\label{sec:ablation-risultati}

Questa sottosezione confronta il modello Pro ($\alpha=0.08$, riferimento)
con i cinque preset di ablation presenti nella Tabella~\ref{tab:ablation-preset}:
\texttt{paper} (replica fedele del riferimento originale), \texttt{environment}
(campo di visibilità, $w_i^t$ e maschera azioni rimossi),
\texttt{temporal} (burn-in e BPTT disattivati), \texttt{single\_dqn}
(Double DQN sostituita da DQN) e \texttt{vanilla\_buffer} (PER, Huber,
gradient clipping, warm-up ed esplorazione ciclica disattivati),
tutti allo stesso $\alpha=0.08$ del riferimento.

\begin{table}[htbp]
  \centering
  \begin{tabular}{@{}lrrr@{}}
    \toprule
    Modello                          & $\overline{tt}$ (s) & $tt_{\max}$ (s) & wait$_{\max}$ (s) \\
    \midrule
    Fixed-Time                       & 339.3               & 1030.0          & 935.0             \\
    Max-Pressure                     & \textbf{161.9}      & 1070.0          & 1040.0            \\
    Pro ($\alpha=0.08$, riferimento) & 172.8               & \textbf{630.0}  & \textbf{425.0}    \\
    \texttt{paper}                   & 187.8               & 950.0           & 915.0             \\
    \texttt{environment}             & 188.3               & 1335.0          & 1290.0            \\
    \texttt{temporal}                & 176.0               & 715.0           & 560.0             \\
    \texttt{single\_dqn}             & 171.9               & \textbf{600.0}  & \textbf{485.0}    \\
    \texttt{vanilla\_buffer}         & \textbf{171.5}      & 695.0           & 560.0             \\
    \bottomrule
  \end{tabular}
  \caption{Travel time medio, massimo e max wait sui preset
    di ablazione, mediati sulle 3 configurazioni di training.}
  \label{tab:ablation-train}
\end{table}

\begin{table}[htbp]
  \centering
  \begin{tabular}{@{}lrrr@{}}
    \toprule
    Modello                          & $\overline{tt}$ (s) & $tt_{\max}$ (s) & wait$_{\max}$ (s) \\
    \midrule
    Fixed-Time                       & 460.7               & 1496.1          & 1311.0            \\
    Max-Pressure                     & \textbf{238.3}      & 1647.4          & 1636.7            \\
    Pro ($\alpha=0.08$, riferimento) & \textbf{241.0}      & \textbf{1017.4} & \textbf{792.0}    \\
    \texttt{paper}                   & 270.4               & 1553.6          & 1513.3            \\
    \texttt{environment}             & 268.6               & 1495.3          & 1463.1            \\
    \texttt{temporal}                & 245.9               & 1129.3          & 940.7             \\
    \texttt{single\_dqn}             & 243.2               & 1051.7          & 917.6             \\
    \texttt{vanilla\_buffer}         & 241.1               & \textbf{1044.9} & \textbf{804.0}    \\
    \bottomrule
  \end{tabular}
  \caption{Come Tabella~\ref{tab:ablation-train}, mediata sulle 7
    configurazioni di test.}
  \label{tab:ablation-test}
\end{table}

Sul travel time medio in training \texttt{vanilla\_buffer},
\texttt{single\_dqn} e il modello Pro restano in
una fascia stretta (171.5-172.8s), con \texttt{temporal} poco distante (176.0s);
\texttt{environment} se ne stacca già in questo regime (188.3s), non solo in test.
In test \texttt{vanilla\_buffer} e Pro restano quasi coincidenti (241.0-241.1s),
seguiti da \texttt{single\_dqn} (243.2s) e \texttt{temporal} (245.9s), mentre
\texttt{environment} e \texttt{paper} restano nettamente distanti (268.6s e 270.4s).
Il distacco di \texttt{environment} non è concentrato sulle due
configurazioni a 200m come si potrebbe attendere dalla sola disattivazione
del cutoff di visibilità: la variazione percentuale rispetto al riferimento
Pro, calcolata per singola configurazione, è di circa il 10\% sulle due
config a 200m e di circa il 12\% sulle restanti cinque config di test (4x4
a 100m, $5\times5$, $6\times6$), un divario distribuito in modo pressoché
uniforme su tutte le configurazioni fuori distribuzione, non specifico
della spaziatura. Questo rende meno convincente l'ipotesi che il solo
cutoff di visibilità sia responsabile del calo: \texttt{environment}
disattiva contemporaneamente anche $w_i^t$ nello stato, la maschera
anti-starvation e il termine di attesa nella ricompensa
(Tabella~\ref{tab:differenze}), e l'effetto osservato è più probabilmente
il risultato congiunto di queste rimozioni che di una sola.

\begin{figure}[htbp]
  \centering
  \begin{subfigure}[t]{0.48\textwidth}
    \centering
    \includegraphics[width=\linewidth]{ablation_train.png}
    \caption{Configurazioni di training}
  \end{subfigure}
  \hfill
  \begin{subfigure}[t]{0.48\textwidth}
    \centering
    \includegraphics[width=\linewidth]{ablation_test.png}
    \caption{Configurazioni di test}
  \end{subfigure}
  \caption{Travel time medio, max wait e travel time massimo sui preset di
    ablation, stessi dati delle
    Tabelle~\ref{tab:ablation-train}-\ref{tab:ablation-test}.}
  \label{fig:overlay-ablation}
\end{figure}

\begin{figure}[htbp]
  \centering
  \begin{subfigure}[t]{0.48\textwidth}
    \centering
    \includegraphics[width=\linewidth]{diff_pct_ablation_train.png}
    \caption{Configurazioni di training}
  \end{subfigure}
  \hfill
  \begin{subfigure}[t]{0.48\textwidth}
    \centering
    \includegraphics[width=\linewidth]{diff_pct_ablation_test.png}
    \caption{Configurazioni di test}
  \end{subfigure}
  \caption{Variazione percentuale di travel time medio e massimo rispetto a
    Max-Pressure (riferimento a 0\%) sui preset di ablation, calcolata per
    singola configurazione/seed e poi mediata.}
  \label{fig:diffpct-ablation}
\end{figure}

Su $tt_{\max}$ e $wait_{max}$, \texttt{paper} ed \texttt{environment} restano
nettamente indietro in entrambi i regimi. I due preset condividono più di
un solo elemento rispetto a ogni altro modello di questo studio: oltre
all'assenza del termine legato all'attesa nella ricompensa, rimuovono
entrambi anche il segnale di attesa $w_i^t$ dallo stato, la maschera
anti-starvation e il vincolo sul campo visivo
(Tabella~\ref{tab:ablation-preset}). Che la sola assenza del termine di
ricompensa non basti a spiegare il divario lo conferma il confronto con
$\alpha=0.00$ (§\ref{sec:risultati-alpha}): pur privo anch'esso di
qualunque penalità sull'attesa nella ricompensa, con stato e maschera
lasciati intatti resta nettamente sotto \texttt{paper} ed
\texttt{environment} su entrambe le metriche di caso peggiore in test
($wait_{\max}$ 1012.3s contro 1513.3s e 1463.1s, $tt_{\max}$ 1164.4s
contro 1553.6s e 1495.3s). Il degrado osservato dipende quindi più
probabilmente dalla combinazione delle rimozioni congiunte che dal solo
termine di ricompensa.

Sempre su queste due metriche, in training \texttt{single\_dqn} ottiene il
miglior $tt_{\max}$ in assoluto (600.0s), con il modello Pro secondo
(630.0s); sul $wait_{\max}$ è invece Pro a ottenere il miglior risultato
del confronto (425.0s), davanti a \texttt{single\_dqn} (485.0s).
\texttt{temporal}, che disattiva solo burn-in e BPTT, è qui il preset con
il $tt_{\max}$ peggiore tra questi quattro (715.0s) e pareggia
\texttt{vanilla\_buffer} sul $wait_{\max}$ (560.0s).

In test \texttt{temporal} resta peggiore di Pro su entrambe le metriche
(1129.3s contro 1017.4s di $tt_{\max}$, 940.7s contro 792.0s di
$wait_{\max}$). Il divario cresce gradualmente con la scala della griglia,
calcolato per singola configurazione rispetto al riferimento Pro: circa
il 10\% sulle tre configurazioni $4\times4$, l'11\% sulla $5\times5$, il
12\% sulla $6\times6$. Una causa plausibile
resta l'assenza di burn-in e BPTT: senza propagazione del gradiente lungo
la sequenza né un periodo di adattamento dello stato ricorrente all'inizio
di ogni episodio, il modulo temporale della Meta-LSTM ha meno modo di
generalizzare il proprio stato interno a condizioni diverse da quelle di
training; questo lavoro non dispone però di un'evidenza indipendente,
oltre all'andamento per configurazione appena descritto, che permetta di
isolare questa causa da altre possibili.

\begin{sloppypar}
In test Pro ottiene il $tt_{\max}$ e il $wait_{\max}$ migliori anche tra
\texttt{single\_dqn} e \texttt{vanilla\_buffer} (1017.4s e 792.0s, contro
1051.7s e 917.6s di \texttt{single\_dqn}, 1044.9s e 804.0s di
\texttt{vanilla\_buffer}): questi due preset non mantengono alcun
vantaggio sistematico sul riferimento.
\end{sloppypar}

\subsection{Studio sui meccanismi spaziali: risultati}
\label{sec:confronto-spaziale}

Questa sottosezione confronta il modello Pro (MetaSTGAT, un layer)
con le sue varianti a meccanismo spaziale alternativo definite in
§\ref{sec:spatial-alt}: MetaSTGAT a due layer, MetaSTGCN a uno e due
layer, e MetaSTSONAR con $L=2$ e $L=4$ ricorrenze. Tutti i modelli
condividono il resto dell'architettura, gli iperparametri e la procedura di training.


\begin{table}[htbp]
  \centering
  \begin{tabular}{@{}lrrr@{}}
    \toprule
    Modello                          & $\overline{tt}$ (s) & $tt_{\max}$ (s) & wait$_{\max}$ (s) \\
    \midrule
    Fixed-Time                       & 339.3               & 1030.0          & 935.0             \\
    Max-Pressure                     & \textbf{161.9}      & 1070.0          & 1040.0            \\
    MetaSTGAT, 1L (Pro, riferimento) & 172.8               & \textbf{630.0}  & \textbf{425.0}    \\
    MetaSTGAT, 2L                    & 181.5               & 685.0           & 525.0             \\
    MetaSTGCN, 1L                    & 260.6               & 1655.0          & 1625.0            \\
    MetaSTGCN, 2L                    & 201.6               & 865.0           & 835.0             \\
    MetaSTSONAR, $L{=}2$             & 179.3               & 710.0           & 575.0             \\
    MetaSTSONAR, $L{=}4$             & \textbf{172.1}      & \textbf{660.0}  & \textbf{475.0}    \\
    \bottomrule
  \end{tabular}
  \caption{Travel time medio, massimo e max wait sui
    meccanismi spaziali, mediati sulle 3 configurazioni di training.}
  \label{tab:architetture-train}
\end{table}

\begin{table}[htbp]
  \centering
  \begin{tabular}{@{}lrrr@{}}
    \toprule
    Modello                          & $\overline{tt}$ (s) & $tt_{\max}$ (s) & wait$_{\max}$ (s) \\
    \midrule
    Fixed-Time                       & 460.7               & 1496.1          & 1311.0            \\
    Max-Pressure                     & \textbf{238.3}      & 1647.4          & 1636.7            \\
    MetaSTGAT, 1L (Pro, riferimento) & \textbf{241.0}      & \textbf{1017.4} & \textbf{792.0}    \\
    MetaSTGAT, 2L                    & 252.7               & \textbf{1037.6} & \textbf{890.6}    \\
    MetaSTGCN, 1L                    & 348.3               & 1749.0          & 1738.7            \\
    MetaSTGCN, 2L                    & 281.1               & 1233.9          & 1170.9            \\
    MetaSTSONAR, $L{=}2$             & 249.8               & 1134.0          & 956.6             \\
    MetaSTSONAR, $L{=}4$             & 246.8               & 1092.9          & 932.6             \\
    \bottomrule
  \end{tabular}
  \caption{Come Tabella~\ref{tab:architetture-train}, mediata sulle 7
    configurazioni di test.}
  \label{tab:architetture-test}
\end{table}

Il divario più netto di questo confronto separa MetaSTGCN dagli altri due
meccanismi. In test, entrambe le varianti
MetaSTGCN registrano valori nettamente peggiori su ogni colonna (travel time medio
281.1-348.3s contro 241.0-252.7s di MetaSTGAT/MetaSTSONAR); MetaSTGCN a un
layer supera persino Max-Pressure in negativo su $tt_{\max}$ e $wait_{\max}$
(1749.0s e 1738.7s, contro 1647.4s e 1636.7s di Max-Pressure), pur
trattandosi di un modello addestrato con lo stesso obiettivo di equità del
riferimento. Anche sul completamento MetaSTGCN a un layer resta la sola
eccezione oltre a Fixed-Time già anticipata in
§\ref{sec:risultati-base} (88.5\% in training, 85.4\% in test, contro il
94.4-97.9\%/89.0-95.8\% di tutte le altre varianti). La differenza
concettuale descritta in §\ref{subsec:meta-gcn},
un peso di aggregazione fissato dalla sola topologia contro un punteggio di
compatibilità appreso per coppia di nodi, si traduce quindi in un divario di
prestazioni ampio su questo compito.
Un secondo layer in MetaSTGCN
migliora la variante a singolo layer su tutte le metriche riportate in
entrambi i regimi ($\overline{tt}$ di training 260.6$\to$201.6s,
$tt_{\max}$ di test 1749.0$\to$1233.9s, $wait_{\max}$ di test
1738.7$\to$1170.9s): la profondità aggiuntiva compensa qui, almeno in
parte, l'assenza di un meccanismo di compatibilità appreso, pur restando
lontana dalle prestazioni di MetaSTGAT e MetaSTSONAR.

MetaSTGAT e MetaSTSONAR hanno prestazioni pressoché pari tra loro su travel
time medio; sul $wait_{\max}$ di test, però, è il modello Pro di
riferimento a ottenere il risultato migliore dell'intero confronto
(792.0s), davanti sia alla propria variante a due layer (890.6s) sia a
entrambe le profondità di MetaSTSONAR (956.6s e 932.6s).

Tra le due profondità di MetaSTSONAR, in questo studio è $L=4$ a risultare
leggermente migliore di $L=2$ sul $tt_{\max}$ in entrambi i regimi (660.0s
contro 710.0s in training, 1092.9s contro 1134.0s in test), coerente con
il vantaggio teorico di campo recettivo più ampio atteso da una profondità
doppia.

Aggiungere un secondo layer a MetaSTGAT, a differenza di MetaSTGCN, peggiora
qui le prestazioni su ogni metrica riportata: travel time medio (172.8s
$\to$ 181.5s in training, 241.0s $\to$ 252.7s in test), $tt_{\max}$
(630.0s $\to$ 685.0s in training, 1017.4s $\to$ 1037.6s in test) e
$wait_{\max}$ (425.0s $\to$ 525.0s in training, 792.0s $\to$ 890.6s in
test, +12\%).

Un caso particolare nelle configurazioni di
training (Tabella~\ref{tab:wait-full}): MetaSTGCN a un layer ottiene un
$wait_{\max}$ medio di 1625s, peggiore di Max-Pressure stesso (1040s), pur
condividendo con il riferimento la stessa ricompensa orientata
all'equità. È il solo modello di tutto questo lavoro, tra quelli
addestrati con il termine anti-starvation pieno, a non riuscire a battere
la baseline reattiva su questa metrica nel proprio stesso regime di
addestramento, coerentemente con la spiegazione architetturale data sopra: un
peso di aggregazione fissato dalla topologia non ha modo di imparare a
dare priorità al segnale del vicino più in sofferenza, per quanto la
ricompensa lo incentivi a farlo.

\begin{figure}[htbp]
  \centering
  \begin{subfigure}[t]{0.48\textwidth}
    \centering
    \includegraphics[width=\linewidth]{architectures_train.png}
    \caption{Configurazioni di training}
  \end{subfigure}
  \hfill
  \begin{subfigure}[t]{0.48\textwidth}
    \centering
    \includegraphics[width=\linewidth]{architectures_test.png}
    \caption{Configurazioni di test}
  \end{subfigure}
  \caption{Travel time medio, max wait e travel time massimo sui meccanismi
    spaziali, stessi dati delle
    Tabelle~\ref{tab:architetture-train}-\ref{tab:architetture-test}.}
  \label{fig:overlay-architetture}
\end{figure}

\begin{figure}[htbp]
  \centering
  \begin{subfigure}[t]{0.48\textwidth}
    \centering
    \includegraphics[width=\linewidth]{diff_pct_architetture_train.png}
    \caption{Configurazioni di training}
  \end{subfigure}
  \hfill
  \begin{subfigure}[t]{0.48\textwidth}
    \centering
    \includegraphics[width=\linewidth]{diff_pct_architetture_test.png}
    \caption{Configurazioni di test}
  \end{subfigure}
  \caption{Variazione percentuale di travel time medio e massimo rispetto a
    Max-Pressure (riferimento a 0\%) sui meccanismi spaziali, calcolata per
    singola configurazione/seed e poi mediata.}
  \label{fig:diffpct-architetture}
\end{figure}

\subsection{Tabelle dei risultati}
Sono qui di seguito riportati i risultati completi di tutti i modelli rispetto a ciascuna configurazione testata.

\begin{sidewaystable}
  \centering
  \tiny
  \begin{tabular}{@{}lrrrrrrrrrr@{}}
    \toprule
    Modello                   & Train 1 & Train 2 & Train 3 & 4x4 Pk & 4x4-200 F & 4x4-200 Pk & 5x5 F & 5x5 Pk & 6x6 F & 6x6 Pk \\
    \midrule
    Fixed-Time                & 84.8    & 85.7    & 88.1    & 85.5   & 85.1      & 82.9       & 86.1  & 83.8   & 83.3  & 81.6   \\
    Max-Pressure              & 98.0    & 98.3    & 97.2    & 92.9   & 93.5      & 92.2       & 92.4  & 91.4   & 91.2  & 89.6   \\
    Pro ($\alpha$=0.08)       & 97.3    & 97.8    & 97.1    & 94.0   & 95.4      & 93.1       & 94.6  & 92.6   & 94.0  & 93.0   \\
    Pro ($\alpha$=0.04)       & 97.7    & 97.7    & 97.1    & 94.0   & 95.8      & 93.2       & 94.4  & 93.5   & 93.7  & 93.2   \\
    Pro ($\alpha$=0.00)       & 96.6    & 97.2    & 97.2    & 93.9   & 95.4      & 93.2       & 94.2  & 92.7   & 93.3  & 91.8   \\
    MetaSTGAT 2L              & 96.5    & 97.3    & 96.9    & 93.2   & 95.3      & 92.1       & 93.7  & 92.6   & 93.2  & 92.5   \\
    Ablation: paper           & 95.0    & 96.6    & 96.3    & 91.8   & 94.4      & 90.8       & 91.9  & 91.2   & 91.4  & 89.4   \\
    Ablation: temporal        & 96.8    & 97.4    & 97.3    & 93.5   & 95.7      & 93.4       & 93.6  & 91.8   & 92.9  & 92.0   \\
    Ablation: single\_dqn     & 97.3    & 97.7    & 97.8    & 93.9   & 95.5      & 92.8       & 94.5  & 92.6   & 93.9  & 90.4   \\
    Ablation: environment     & 95.4    & 96.1    & 96.5    & 92.1   & 93.4      & 91.2       & 92.0  & 89.5   & 91.3  & 89.0   \\
    Ablation: vanilla\_buffer & 96.1    & 97.7    & 97.5    & 93.9   & 95.7      & 93.3       & 94.2  & 93.1   & 93.4  & 92.4   \\
    MetaSTGCN 1L              & 89.1    & 88.5    & 87.9    & 85.7   & 84.2      & 82.8       & 87.4  & 86.4   & 87.6  & 83.9   \\
    MetaSTGCN 2L              & 95.0    & 94.4    & 96.3    & 91.4   & 93.7      & 91.1       & 92.2  & 90.6   & 91.7  & 90.3   \\
    MetaSTSONAR L=2           & 96.9    & 97.8    & 96.8    & 93.9   & 95.7      & 93.5       & 93.8  & 92.8   & 93.5  & 92.3   \\
    MetaSTSONAR L=4           & 97.3    & 97.7    & 97.9    & 93.9   & 95.4      & 93.4       & 93.9  & 92.3   & 93.0  & 92.0   \\
    \bottomrule
  \end{tabular}
  \caption{Percentuale di trip completion $tc$ per ogni modello
    e ogni configurazione, corretta escludendo dal
    denominatore i veicoli strutturalmente impossibilitati ad arrivare in
    tempo. Nessun modello scende sotto l'80\%; a parte Fixed-Time (baseline non
    adattiva) l'unico caso ricorrente sotto il 90\% è MetaSTGCN 1L.}
  \label{tab:tc-adjusted}
\end{sidewaystable}

\begin{sidewaystable}
  \centering
  \tiny
  \begin{tabular}{@{}lrrrrrrrrrr@{}}
    \toprule
    Modello                   & Train 1 & Train 2 & Train 3 & 4x4 Pk                                     & 4x4-200 F                                   & 4x4-200 Pk                                 & 5x5 F                                       & 5x5 Pk                                     & 6x6 F                                       & 6x6 Pk                                      \\
    \midrule
    Fixed-Time                & 351.6   & 342.9   & 323.3   & 437.3 {\fontsize{4}{4}\selectfont$\pm$6.4} & 415.5 {\fontsize{4}{4}\selectfont$\pm$6.7}  & 438.8 {\fontsize{4}{4}\selectfont$\pm$4.9} & 459.7 {\fontsize{4}{4}\selectfont$\pm$4.7}  & 485.0 {\fontsize{4}{4}\selectfont$\pm$5.1} & 484.8 {\fontsize{4}{4}\selectfont$\pm$6.2}  & 503.7 {\fontsize{4}{4}\selectfont$\pm$4.6}  \\
    Max-Pressure              & 165.4   & 160.2   & 160.0   & 218.8 {\fontsize{4}{4}\selectfont$\pm$2.3} & 197.3 {\fontsize{4}{4}\selectfont$\pm$4.2}  & 247.8 {\fontsize{4}{4}\selectfont$\pm$2.6} & 217.8 {\fontsize{4}{4}\selectfont$\pm$3.3}  & 258.6 {\fontsize{4}{4}\selectfont$\pm$4.9} & 244.7 {\fontsize{4}{4}\selectfont$\pm$5.5}  & 283.3 {\fontsize{4}{4}\selectfont$\pm$1.4}  \\
    Pro ($\alpha$=0.08)       & 178.0   & 174.8   & 165.5   & 229.4 {\fontsize{4}{4}\selectfont$\pm$4.1} & 202.8 {\fontsize{4}{4}\selectfont$\pm$3.2}  & 253.3 {\fontsize{4}{4}\selectfont$\pm$3.4} & 220.7 {\fontsize{4}{4}\selectfont$\pm$5.0}  & 268.0 {\fontsize{4}{4}\selectfont$\pm$7.5} & 238.1 {\fontsize{4}{4}\selectfont$\pm$4.2}  & 275.1 {\fontsize{4}{4}\selectfont$\pm$3.4}  \\
    Pro ($\alpha$=0.04)       & 175.1   & 171.7   & 165.2   & 229.3 {\fontsize{4}{4}\selectfont$\pm$2.8} & 202.4 {\fontsize{4}{4}\selectfont$\pm$2.6}  & 251.7 {\fontsize{4}{4}\selectfont$\pm$2.5} & 221.4 {\fontsize{4}{4}\selectfont$\pm$4.8}  & 264.1 {\fontsize{4}{4}\selectfont$\pm$3.0} & 242.6 {\fontsize{4}{4}\selectfont$\pm$2.6}  & 276.2 {\fontsize{4}{4}\selectfont$\pm$3.2}  \\
    Pro ($\alpha$=0.00)       & 177.8   & 169.5   & 173.1   & 230.8 {\fontsize{4}{4}\selectfont$\pm$2.7} & 201.3 {\fontsize{4}{4}\selectfont$\pm$6.3}  & 250.5 {\fontsize{4}{4}\selectfont$\pm$3.5} & 226.5 {\fontsize{4}{4}\selectfont$\pm$6.5}  & 275.0 {\fontsize{4}{4}\selectfont$\pm$6.2} & 253.4 {\fontsize{4}{4}\selectfont$\pm$4.0}  & 295.0 {\fontsize{4}{4}\selectfont$\pm$7.8}  \\
    MetaSTGAT 2L              & 188.8   & 177.6   & 178.1   & 244.8 {\fontsize{4}{4}\selectfont$\pm$3.2} & 208.3 {\fontsize{4}{4}\selectfont$\pm$3.6}  & 262.6 {\fontsize{4}{4}\selectfont$\pm$3.4} & 234.6 {\fontsize{4}{4}\selectfont$\pm$4.9}  & 279.1 {\fontsize{4}{4}\selectfont$\pm$1.7} & 247.4 {\fontsize{4}{4}\selectfont$\pm$5.8}  & 291.9 {\fontsize{4}{4}\selectfont$\pm$6.4}  \\
    Ablation: paper           & 194.1   & 182.9   & 186.5   & 253.9 {\fontsize{4}{4}\selectfont$\pm$4.6} & 217.7 {\fontsize{4}{4}\selectfont$\pm$3.6}  & 287.3 {\fontsize{4}{4}\selectfont$\pm$4.8} & 245.3 {\fontsize{4}{4}\selectfont$\pm$11.4} & 294.5 {\fontsize{4}{4}\selectfont$\pm$5.2} & 275.8 {\fontsize{4}{4}\selectfont$\pm$2.6}  & 318.6 {\fontsize{4}{4}\selectfont$\pm$3.5}  \\
    Ablation: temporal        & 181.7   & 173.0   & 173.4   & 234.4 {\fontsize{4}{4}\selectfont$\pm$3.9} & 203.7 {\fontsize{4}{4}\selectfont$\pm$4.9}  & 253.0 {\fontsize{4}{4}\selectfont$\pm$3.4} & 224.8 {\fontsize{4}{4}\selectfont$\pm$7.1}  & 272.9 {\fontsize{4}{4}\selectfont$\pm$4.5} & 246.3 {\fontsize{4}{4}\selectfont$\pm$5.0}  & 285.8 {\fontsize{4}{4}\selectfont$\pm$5.6}  \\
    Ablation: single\_dqn     & 179.0   & 172.9   & 163.8   & 228.1 {\fontsize{4}{4}\selectfont$\pm$3.6} & 201.1 {\fontsize{4}{4}\selectfont$\pm$4.4}  & 253.3 {\fontsize{4}{4}\selectfont$\pm$2.8} & 219.9 {\fontsize{4}{4}\selectfont$\pm$3.2}  & 270.2 {\fontsize{4}{4}\selectfont$\pm$8.1} & 237.7 {\fontsize{4}{4}\selectfont$\pm$4.9}  & 291.9 {\fontsize{4}{4}\selectfont$\pm$7.3}  \\
    Ablation: environment     & 195.8   & 186.8   & 182.4   & 253.5 {\fontsize{4}{4}\selectfont$\pm$6.9} & 222.9 {\fontsize{4}{4}\selectfont$\pm$4.5}  & 279.1 {\fontsize{4}{4}\selectfont$\pm$3.9} & 242.2 {\fontsize{4}{4}\selectfont$\pm$7.6}  & 291.9 {\fontsize{4}{4}\selectfont$\pm$2.6} & 274.1 {\fontsize{4}{4}\selectfont$\pm$8.4}  & 316.5 {\fontsize{4}{4}\selectfont$\pm$7.8}  \\
    Ablation: vanilla\_buffer & 177.9   & 170.2   & 166.5   & 230.3 {\fontsize{4}{4}\selectfont$\pm$2.9} & 202.2 {\fontsize{4}{4}\selectfont$\pm$4.1}  & 252.4 {\fontsize{4}{4}\selectfont$\pm$4.2} & 220.1 {\fontsize{4}{4}\selectfont$\pm$4.5}  & 264.8 {\fontsize{4}{4}\selectfont$\pm$3.9} & 238.8 {\fontsize{4}{4}\selectfont$\pm$3.2}  & 278.9 {\fontsize{4}{4}\selectfont$\pm$6.4}  \\
    MetaSTGCN 1L              & 257.4   & 262.4   & 261.9   & 335.7 {\fontsize{4}{4}\selectfont$\pm$9.4} & 328.4 {\fontsize{4}{4}\selectfont$\pm$15.3} & 360.2 {\fontsize{4}{4}\selectfont$\pm$5.5} & 323.4 {\fontsize{4}{4}\selectfont$\pm$22.7} & 361.6 {\fontsize{4}{4}\selectfont$\pm$3.7} & 331.9 {\fontsize{4}{4}\selectfont$\pm$12.2} & 397.2 {\fontsize{4}{4}\selectfont$\pm$13.9} \\
    MetaSTGCN 2L              & 206.7   & 205.9   & 192.2   & 277.0 {\fontsize{4}{4}\selectfont$\pm$6.4} & 228.6 {\fontsize{4}{4}\selectfont$\pm$6.1}  & 280.6 {\fontsize{4}{4}\selectfont$\pm$3.9} & 259.9 {\fontsize{4}{4}\selectfont$\pm$4.3}  & 317.8 {\fontsize{4}{4}\selectfont$\pm$2.0} & 279.6 {\fontsize{4}{4}\selectfont$\pm$4.5}  & 324.3 {\fontsize{4}{4}\selectfont$\pm$5.9}  \\
    MetaSTSONAR L=2           & 181.3   & 174.4   & 182.2   & 234.6 {\fontsize{4}{4}\selectfont$\pm$3.4} & 206.1 {\fontsize{4}{4}\selectfont$\pm$5.4}  & 258.3 {\fontsize{4}{4}\selectfont$\pm$2.4} & 229.0 {\fontsize{4}{4}\selectfont$\pm$3.9}  & 276.9 {\fontsize{4}{4}\selectfont$\pm$2.8} & 251.1 {\fontsize{4}{4}\selectfont$\pm$4.6}  & 292.9 {\fontsize{4}{4}\selectfont$\pm$7.5}  \\
    MetaSTSONAR L=4           & 176.5   & 170.5   & 169.3   & 232.2 {\fontsize{4}{4}\selectfont$\pm$3.9} & 204.4 {\fontsize{4}{4}\selectfont$\pm$7.6}  & 252.0 {\fontsize{4}{4}\selectfont$\pm$3.9} & 227.1 {\fontsize{4}{4}\selectfont$\pm$5.1}  & 273.0 {\fontsize{4}{4}\selectfont$\pm$3.6} & 248.9 {\fontsize{4}{4}\selectfont$\pm$3.5}  & 290.2 {\fontsize{4}{4}\selectfont$\pm$1.6}  \\
    \bottomrule
  \end{tabular}
  \caption{Travel time medio $\overline{tt}$ (s) per ogni modello e ogni configurazione. Sulle configurazioni di test, media $\pm$ deviazione standard su 5 seed diversi.}
  \label{tab:ttmedio-full}
\end{sidewaystable}

\begin{sidewaystable}
  \centering
  \tiny
  \begin{tabular}{@{}lrrrrrrrrrr@{}}
    \toprule
    Modello                   & Train 1 & Train 2 & Train 3 & 4x4 Pk                                        & 4x4-200 F                                     & 4x4-200 Pk                                    & 5x5 F                                         & 5x5 Pk                                        & 6x6 F                                         & 6x6 Pk                                        \\
    \midrule
    Fixed-Time                & 1020.0  & 975.0   & 1095.0  & 1410.0 {\fontsize{4}{4}\selectfont$\pm$49.3}  & 1431.0 {\fontsize{4}{4}\selectfont$\pm$30.9}  & 1290.0 {\fontsize{4}{4}\selectfont$\pm$19.0}  & 1620.0 {\fontsize{4}{4}\selectfont$\pm$32.9}  & 1482.0 {\fontsize{4}{4}\selectfont$\pm$17.5}  & 1653.0 {\fontsize{4}{4}\selectfont$\pm$32.0}  & 1587.0 {\fontsize{4}{4}\selectfont$\pm$37.2}  \\
    Max-Pressure              & 915.0   & 1155.0  & 1140.0  & 1194.0 {\fontsize{4}{4}\selectfont$\pm$20.3}  & 1692.0 {\fontsize{4}{4}\selectfont$\pm$133.6} & 1446.0 {\fontsize{4}{4}\selectfont$\pm$105.0} & 1800.0 {\fontsize{4}{4}\selectfont$\pm$0.0}   & 1800.0 {\fontsize{4}{4}\selectfont$\pm$0.0}   & 1800.0 {\fontsize{4}{4}\selectfont$\pm$0.0}   & 1800.0 {\fontsize{4}{4}\selectfont$\pm$0.0}   \\
    Pro ($\alpha$=0.08)       & 600.0   & 600.0   & 690.0   & 861.0 {\fontsize{4}{4}\selectfont$\pm$45.1}   & 774.0 {\fontsize{4}{4}\selectfont$\pm$55.8}   & 918.0 {\fontsize{4}{4}\selectfont$\pm$51.4}   & 1074.0 {\fontsize{4}{4}\selectfont$\pm$113.7} & 1083.0 {\fontsize{4}{4}\selectfont$\pm$66.7}  & 1269.0 {\fontsize{4}{4}\selectfont$\pm$87.3}  & 1143.0 {\fontsize{4}{4}\selectfont$\pm$36.0}  \\
    Pro ($\alpha$=0.04)       & 660.0   & 705.0   & 690.0   & 924.0 {\fontsize{4}{4}\selectfont$\pm$48.0}   & 819.0 {\fontsize{4}{4}\selectfont$\pm$33.7}   & 1002.0 {\fontsize{4}{4}\selectfont$\pm$112.8} & 1092.0 {\fontsize{4}{4}\selectfont$\pm$53.2}  & 1092.0 {\fontsize{4}{4}\selectfont$\pm$49.7}  & 1260.0 {\fontsize{4}{4}\selectfont$\pm$28.5}  & 1194.0 {\fontsize{4}{4}\selectfont$\pm$20.3}  \\
    Pro ($\alpha$=0.00)       & 720.0   & 645.0   & 855.0   & 1044.0 {\fontsize{4}{4}\selectfont$\pm$30.9}  & 840.0 {\fontsize{4}{4}\selectfont$\pm$47.4}   & 954.0 {\fontsize{4}{4}\selectfont$\pm$12.0}   & 1269.0 {\fontsize{4}{4}\selectfont$\pm$94.2}  & 1149.0 {\fontsize{4}{4}\selectfont$\pm$56.6}  & 1482.0 {\fontsize{4}{4}\selectfont$\pm$163.1} & 1413.0 {\fontsize{4}{4}\selectfont$\pm$207.9} \\
    MetaSTGAT 2L              & 660.0   & 600.0   & 795.0   & 915.0 {\fontsize{4}{4}\selectfont$\pm$44.5}   & 786.0 {\fontsize{4}{4}\selectfont$\pm$48.9}   & 930.0 {\fontsize{4}{4}\selectfont$\pm$26.8}   & 1161.0 {\fontsize{4}{4}\selectfont$\pm$51.6}  & 1101.0 {\fontsize{4}{4}\selectfont$\pm$51.6}  & 1224.0 {\fontsize{4}{4}\selectfont$\pm$68.8}  & 1146.0 {\fontsize{4}{4}\selectfont$\pm$35.0}  \\
    Ablation: paper           & 930.0   & 975.0   & 945.0   & 1212.0 {\fontsize{4}{4}\selectfont$\pm$142.7} & 1257.0 {\fontsize{4}{4}\selectfont$\pm$93.1}  & 1326.0 {\fontsize{4}{4}\selectfont$\pm$145.3} & 1800.0 {\fontsize{4}{4}\selectfont$\pm$0.0}   & 1794.0 {\fontsize{4}{4}\selectfont$\pm$12.0}  & 1800.0 {\fontsize{4}{4}\selectfont$\pm$0.0}   & 1686.0 {\fontsize{4}{4}\selectfont$\pm$152.0} \\
    Ablation: temporal        & 735.0   & 585.0   & 825.0   & 987.0 {\fontsize{4}{4}\selectfont$\pm$48.7}   & 861.0 {\fontsize{4}{4}\selectfont$\pm$72.6}   & 963.0 {\fontsize{4}{4}\selectfont$\pm$44.9}   & 1191.0 {\fontsize{4}{4}\selectfont$\pm$95.6}  & 1203.0 {\fontsize{4}{4}\selectfont$\pm$63.9}  & 1470.0 {\fontsize{4}{4}\selectfont$\pm$112.2} & 1230.0 {\fontsize{4}{4}\selectfont$\pm$37.9}  \\
    Ablation: single\_dqn     & 600.0   & 585.0   & 615.0   & 825.0 {\fontsize{4}{4}\selectfont$\pm$34.2}   & 759.0 {\fontsize{4}{4}\selectfont$\pm$47.1}   & 948.0 {\fontsize{4}{4}\selectfont$\pm$133.0}  & 1092.0 {\fontsize{4}{4}\selectfont$\pm$128.5} & 1122.0 {\fontsize{4}{4}\selectfont$\pm$65.3}  & 1251.0 {\fontsize{4}{4}\selectfont$\pm$45.1}  & 1365.0 {\fontsize{4}{4}\selectfont$\pm$148.2} \\
    Ablation: environment     & 1425.0  & 1080.0  & 1500.0  & 1239.0 {\fontsize{4}{4}\selectfont$\pm$173.0} & 1440.0 {\fontsize{4}{4}\selectfont$\pm$153.6} & 1296.0 {\fontsize{4}{4}\selectfont$\pm$54.2}  & 1638.0 {\fontsize{4}{4}\selectfont$\pm$161.4} & 1605.0 {\fontsize{4}{4}\selectfont$\pm$103.5} & 1599.0 {\fontsize{4}{4}\selectfont$\pm$185.1} & 1650.0 {\fontsize{4}{4}\selectfont$\pm$195.3} \\
    Ablation: vanilla\_buffer & 660.0   & 675.0   & 750.0   & 954.0 {\fontsize{4}{4}\selectfont$\pm$38.7}   & 795.0 {\fontsize{4}{4}\selectfont$\pm$42.4}   & 936.0 {\fontsize{4}{4}\selectfont$\pm$43.1}   & 1161.0 {\fontsize{4}{4}\selectfont$\pm$102.4} & 1071.0 {\fontsize{4}{4}\selectfont$\pm$33.7}  & 1215.0 {\fontsize{4}{4}\selectfont$\pm$26.8}  & 1182.0 {\fontsize{4}{4}\selectfont$\pm$22.0}  \\
    MetaSTGCN 1L              & 1740.0  & 1425.0  & 1800.0  & 1776.0 {\fontsize{4}{4}\selectfont$\pm$48.0}  & 1791.0 {\fontsize{4}{4}\selectfont$\pm$18.0}  & 1626.0 {\fontsize{4}{4}\selectfont$\pm$179.4} & 1800.0 {\fontsize{4}{4}\selectfont$\pm$0.0}   & 1743.0 {\fontsize{4}{4}\selectfont$\pm$114.0} & 1791.0 {\fontsize{4}{4}\selectfont$\pm$18.0}  & 1716.0 {\fontsize{4}{4}\selectfont$\pm$104.6} \\
    MetaSTGCN 2L              & 930.0   & 855.0   & 810.0   & 1125.0 {\fontsize{4}{4}\selectfont$\pm$82.7}  & 1011.0 {\fontsize{4}{4}\selectfont$\pm$144.4} & 978.0 {\fontsize{4}{4}\selectfont$\pm$29.1}   & 1512.0 {\fontsize{4}{4}\selectfont$\pm$204.0} & 1275.0 {\fontsize{4}{4}\selectfont$\pm$124.1} & 1458.0 {\fontsize{4}{4}\selectfont$\pm$49.7}  & 1278.0 {\fontsize{4}{4}\selectfont$\pm$74.3}  \\
    MetaSTSONAR L=2           & 705.0   & 600.0   & 825.0   & 957.0 {\fontsize{4}{4}\selectfont$\pm$25.8}   & 873.0 {\fontsize{4}{4}\selectfont$\pm$77.3}   & 1047.0 {\fontsize{4}{4}\selectfont$\pm$29.1}  & 1392.0 {\fontsize{4}{4}\selectfont$\pm$199.5} & 1128.0 {\fontsize{4}{4}\selectfont$\pm$42.8}  & 1314.0 {\fontsize{4}{4}\selectfont$\pm$38.7}  & 1227.0 {\fontsize{4}{4}\selectfont$\pm$69.3}  \\
    MetaSTSONAR L=4           & 585.0   & 570.0   & 825.0   & 933.0 {\fontsize{4}{4}\selectfont$\pm$43.9}   & 846.0 {\fontsize{4}{4}\selectfont$\pm$123.9}  & 927.0 {\fontsize{4}{4}\selectfont$\pm$34.7}   & 1287.0 {\fontsize{4}{4}\selectfont$\pm$99.2}  & 1134.0 {\fontsize{4}{4}\selectfont$\pm$54.2}  & 1320.0 {\fontsize{4}{4}\selectfont$\pm$51.1}  & 1203.0 {\fontsize{4}{4}\selectfont$\pm$6.0}   \\
    \bottomrule
  \end{tabular}
  \caption{Travel time massimo $tt_{\max}$ (s) per ogni modello e ogni configurazione. Sulle configurazioni di test, media $\pm$ deviazione standard su 5 seed diversi.}
  \label{tab:ttmax-full}
\end{sidewaystable}

\begin{sidewaystable}
  \centering
  \tiny
  \begin{tabular}{@{}lrrrrrrrrrr@{}}
    \toprule
    Modello                   & Train 1 & Train 2 & Train 3 & 4x4 Pk                                        & 4x4-200 F                                     & 4x4-200 Pk                                    & 5x5 F                                         & 5x5 Pk                                        & 6x6 F                                         & 6x6 Pk                                        \\
    \midrule
    Fixed-Time                & 930.0   & 900.0   & 975.0   & 1047.0 {\fontsize{4}{4}\selectfont$\pm$100.6} & 1053.0 {\fontsize{4}{4}\selectfont$\pm$75.5}  & 1014.0 {\fontsize{4}{4}\selectfont$\pm$64.8}  & 1602.0 {\fontsize{4}{4}\selectfont$\pm$36.0}  & 1335.0 {\fontsize{4}{4}\selectfont$\pm$74.7}  & 1629.0 {\fontsize{4}{4}\selectfont$\pm$48.9}  & 1497.0 {\fontsize{4}{4}\selectfont$\pm$66.7}  \\
    Max-Pressure              & 885.0   & 1125.0  & 1110.0  & 1164.0 {\fontsize{4}{4}\selectfont$\pm$37.5}  & 1680.0 {\fontsize{4}{4}\selectfont$\pm$138.1} & 1419.0 {\fontsize{4}{4}\selectfont$\pm$98.4}  & 1800.0 {\fontsize{4}{4}\selectfont$\pm$0.0}   & 1794.0 {\fontsize{4}{4}\selectfont$\pm$12.0}  & 1800.0 {\fontsize{4}{4}\selectfont$\pm$0.0}   & 1800.0 {\fontsize{4}{4}\selectfont$\pm$0.0}   \\
    Pro ($\alpha$=0.08)       & 465.0   & 360.0   & 450.0   & 591.0 {\fontsize{4}{4}\selectfont$\pm$39.8}   & 549.0 {\fontsize{4}{4}\selectfont$\pm$36.2}   & 756.0 {\fontsize{4}{4}\selectfont$\pm$105.0}  & 726.0 {\fontsize{4}{4}\selectfont$\pm$117.2}  & 933.0 {\fontsize{4}{4}\selectfont$\pm$149.8}  & 948.0 {\fontsize{4}{4}\selectfont$\pm$164.2}  & 1041.0 {\fontsize{4}{4}\selectfont$\pm$77.4}  \\
    Pro ($\alpha$=0.04)       & 510.0   & 630.0   & 540.0   & 846.0 {\fontsize{4}{4}\selectfont$\pm$45.1}   & 630.0 {\fontsize{4}{4}\selectfont$\pm$53.7}   & 885.0 {\fontsize{4}{4}\selectfont$\pm$157.9}  & 765.0 {\fontsize{4}{4}\selectfont$\pm$83.2}   & 1026.0 {\fontsize{4}{4}\selectfont$\pm$97.5}  & 1182.0 {\fontsize{4}{4}\selectfont$\pm$102.4} & 1182.0 {\fontsize{4}{4}\selectfont$\pm$36.0}  \\
    Pro ($\alpha$=0.00)       & 525.0   & 435.0   & 705.0   & 888.0 {\fontsize{4}{4}\selectfont$\pm$95.1}   & 690.0 {\fontsize{4}{4}\selectfont$\pm$45.5}   & 813.0 {\fontsize{4}{4}\selectfont$\pm$59.5}   & 951.0 {\fontsize{4}{4}\selectfont$\pm$125.0}  & 1062.0 {\fontsize{4}{4}\selectfont$\pm$71.9}  & 1320.0 {\fontsize{4}{4}\selectfont$\pm$212.8} & 1362.0 {\fontsize{4}{4}\selectfont$\pm$202.0} \\
    MetaSTGAT 2L              & 480.0   & 495.0   & 600.0   & 714.0 {\fontsize{4}{4}\selectfont$\pm$61.2}   & 564.0 {\fontsize{4}{4}\selectfont$\pm$43.1}   & 792.0 {\fontsize{4}{4}\selectfont$\pm$92.2}   & 975.0 {\fontsize{4}{4}\selectfont$\pm$209.4}  & 1017.0 {\fontsize{4}{4}\selectfont$\pm$121.3} & 1077.0 {\fontsize{4}{4}\selectfont$\pm$179.6} & 1095.0 {\fontsize{4}{4}\selectfont$\pm$74.1}  \\
    Ablation: paper           & 885.0   & 945.0   & 915.0   & 1071.0 {\fontsize{4}{4}\selectfont$\pm$109.7} & 1215.0 {\fontsize{4}{4}\selectfont$\pm$91.5}  & 1278.0 {\fontsize{4}{4}\selectfont$\pm$135.3} & 1800.0 {\fontsize{4}{4}\selectfont$\pm$0.0}   & 1788.0 {\fontsize{4}{4}\selectfont$\pm$24.0}  & 1800.0 {\fontsize{4}{4}\selectfont$\pm$0.0}   & 1641.0 {\fontsize{4}{4}\selectfont$\pm$171.4} \\
    Ablation: temporal        & 555.0   & 510.0   & 615.0   & 786.0 {\fontsize{4}{4}\selectfont$\pm$98.4}   & 609.0 {\fontsize{4}{4}\selectfont$\pm$110.9}  & 783.0 {\fontsize{4}{4}\selectfont$\pm$85.6}   & 846.0 {\fontsize{4}{4}\selectfont$\pm$73.9}   & 1089.0 {\fontsize{4}{4}\selectfont$\pm$121.7} & 1326.0 {\fontsize{4}{4}\selectfont$\pm$155.5} & 1146.0 {\fontsize{4}{4}\selectfont$\pm$51.6}  \\
    Ablation: single\_dqn     & 525.0   & 435.0   & 495.0   & 696.0 {\fontsize{4}{4}\selectfont$\pm$56.6}   & 606.0 {\fontsize{4}{4}\selectfont$\pm$43.1}   & 807.0 {\fontsize{4}{4}\selectfont$\pm$185.0}  & 897.0 {\fontsize{4}{4}\selectfont$\pm$151.9}  & 1047.0 {\fontsize{4}{4}\selectfont$\pm$78.5}  & 1020.0 {\fontsize{4}{4}\selectfont$\pm$100.8} & 1350.0 {\fontsize{4}{4}\selectfont$\pm$159.9} \\
    Ablation: environment     & 1395.0  & 1050.0  & 1425.0  & 1179.0 {\fontsize{4}{4}\selectfont$\pm$190.8} & 1413.0 {\fontsize{4}{4}\selectfont$\pm$163.1} & 1260.0 {\fontsize{4}{4}\selectfont$\pm$50.2}  & 1614.0 {\fontsize{4}{4}\selectfont$\pm$166.7} & 1581.0 {\fontsize{4}{4}\selectfont$\pm$118.7} & 1557.0 {\fontsize{4}{4}\selectfont$\pm$219.5} & 1638.0 {\fontsize{4}{4}\selectfont$\pm$201.3} \\
    Ablation: vanilla\_buffer & 585.0   & 570.0   & 525.0   & 753.0 {\fontsize{4}{4}\selectfont$\pm$85.1}   & 594.0 {\fontsize{4}{4}\selectfont$\pm$50.7}   & 777.0 {\fontsize{4}{4}\selectfont$\pm$59.5}   & 768.0 {\fontsize{4}{4}\selectfont$\pm$235.3}  & 837.0 {\fontsize{4}{4}\selectfont$\pm$17.5}   & 834.0 {\fontsize{4}{4}\selectfont$\pm$99.8}   & 1065.0 {\fontsize{4}{4}\selectfont$\pm$70.4}  \\
    MetaSTGCN 1L              & 1710.0  & 1395.0  & 1770.0  & 1752.0 {\fontsize{4}{4}\selectfont$\pm$59.5}  & 1776.0 {\fontsize{4}{4}\selectfont$\pm$40.9}  & 1602.0 {\fontsize{4}{4}\selectfont$\pm$171.9} & 1800.0 {\fontsize{4}{4}\selectfont$\pm$0.0}   & 1743.0 {\fontsize{4}{4}\selectfont$\pm$114.0} & 1785.0 {\fontsize{4}{4}\selectfont$\pm$30.0}  & 1713.0 {\fontsize{4}{4}\selectfont$\pm$109.2} \\
    MetaSTGCN 2L              & 900.0   & 825.0   & 780.0   & 1086.0 {\fontsize{4}{4}\selectfont$\pm$122.4} & 861.0 {\fontsize{4}{4}\selectfont$\pm$192.9}  & 897.0 {\fontsize{4}{4}\selectfont$\pm$77.9}   & 1443.0 {\fontsize{4}{4}\selectfont$\pm$244.9} & 1233.0 {\fontsize{4}{4}\selectfont$\pm$127.4} & 1419.0 {\fontsize{4}{4}\selectfont$\pm$72.6}  & 1257.0 {\fontsize{4}{4}\selectfont$\pm$74.3}  \\
    MetaSTSONAR L=2           & 600.0   & 420.0   & 705.0   & 771.0 {\fontsize{4}{4}\selectfont$\pm$78.6}   & 645.0 {\fontsize{4}{4}\selectfont$\pm$93.4}   & 948.0 {\fontsize{4}{4}\selectfont$\pm$32.0}   & 1098.0 {\fontsize{4}{4}\selectfont$\pm$328.8} & 975.0 {\fontsize{4}{4}\selectfont$\pm$99.9}   & 1110.0 {\fontsize{4}{4}\selectfont$\pm$152.1} & 1149.0 {\fontsize{4}{4}\selectfont$\pm$60.4}  \\
    MetaSTSONAR L=4           & 480.0   & 435.0   & 510.0   & 786.0 {\fontsize{4}{4}\selectfont$\pm$118.7}  & 609.0 {\fontsize{4}{4}\selectfont$\pm$123.9}  & 807.0 {\fontsize{4}{4}\selectfont$\pm$91.2}   & 1101.0 {\fontsize{4}{4}\selectfont$\pm$240.9} & 1116.0 {\fontsize{4}{4}\selectfont$\pm$63.4}  & 930.0 {\fontsize{4}{4}\selectfont$\pm$164.9}  & 1179.0 {\fontsize{4}{4}\selectfont$\pm$49.8}  \\
    \bottomrule
  \end{tabular}
  \caption{Max wait $\text{wait}_{\max}$ (s) per ogni modello e ogni configurazione. Sulle configurazioni di test, media $\pm$ deviazione standard su 5 seed diversi. Sulle config \texttt{4x4-200}, valori a piena corsia.}
  \label{tab:wait-full}
\end{sidewaystable}
